\documentclass[10pt,a4paper]{amsart}
\usepackage[margin=19mm]{geometry}
\usepackage{amsmath,amssymb,mathtools}
\usepackage[scr=rsfs,cal=euler]{mathalfa}
\usepackage{xfrac}
\usepackage[unicode,hidelinks,bookmarks=false]{hyperref}
\newcommand{\ie}{i.e.\ }
\newcommand{\cf}{cf.\ }
\newcommand{\resp}{respectively\ }
\newcommand{\Subsection}{\subsection}
\providecommand{\nobreakdash}{\nobreak\hbox{-}}
\setlength{\emergencystretch}{2em}
\multlinegap0pt
\usepackage{bm}
\usepackage{bbm}
\usepackage{dsfont}
\usepackage{xspace}
\usepackage{cancel}
\usepackage{mathtools}
\usepackage{amsthm}
\makeatletter
\@ifundefined{theorem}{\newtheorem{theorem}{Theorem}}{}
\@ifundefined{lemma}{\newtheorem{lemma}[theorem]{Lemma}}{}
\makeatother
\title[Applications of Nonlinear Projections to Rectifiable 1-sets]{Applications of Nonlinear Projections to Rectifiable 1-sets}
\author{Rosemarie Bongers, Paige Bright, Caleb Marshall, Krystal Taylor}
\date{Source: arXiv:2608.10253v1 (CC BY 4.0)}
\begin{document}
\maketitle

\noindent\emph{Source and licence.} Applications of Nonlinear Projections to Rectifiable 1-sets. Version arXiv:2608.10253v1 (CC BY 4.0), distributed under the Creative Commons Attribution 4.0 International licence, as recorded in the arXiv licence metadata for this exact version and shown on the abstract page https://arxiv.org/abs/2608.10253v1. Full rights notice: licenses/arxiv-2608.10253-CC-BY-4.0.txt.\par\smallskip

\noindent\textit{Editorial scope.} Counted unit: the Lemma whose source label is \texttt{lemma:circles\_canonical\_embedding} in the article (source0.tex lines 583--585, source label \texttt{lemma:circles\_canonical\_embedding}), delivered together with the article's own complete original proof of it (source0.tex lines 587--621, one \texttt{proof} environment of 35 lines). No mathematics is written, repaired or re-derived: statement and proof are the article's own text line for line. Required context retained uncounted: none. Every definition, display and label of those lines is retained unchanged. Omitted from this package and not redistributed: 1--582, 586--586, 622--654. Nothing inside a retained excerpt is omitted or altered: every retained body is delivered exactly as the article's lines are written, so no omission receipt of any kind accompanies this package. Citations: the retained excerpts carry no citation command at all, so no bibliography entry of the article is transcribed and no reference is redistributed. Reference adaptation: none was needed and none was made; every reference inside the delivered unit resolves inside it, so no unresolved reference or citation is produced. Printed slips kept literal: no printed word, symbol, hypothesis or number of the counted statement or of its proof was altered. Layout and class: the article's own class is \texttt{amsart} with packages including geometry, graphicx, comment, amsmath, amssymb, amsfonts, amsthm, mathrsfs, enumerate, float, fancyhdr, xcolor, hyperref, resmes; the wrapper is the lane's pinned \texttt{amsart} editorial wrapper at 10pt on A4, which restates the theorem-like environments the retained text needs and adds the article's own macro definitions that the retained text needs (none), with extra wrapper packages (bm, bbm, dsfont, xspace, cancel, mathtools). The delivered numbering is the unit's own. Assets: no retained span contains a figure environment, a graphics-inclusion command, a PDF-page inclusion, a table or a TikZ picture; no image file is copied into this package, the package declares no asset dependency and no image file is redistributed with it. Licence: version arXiv:2608.10253v1 of this article is distributed under the Creative Commons Attribution 4.0 International licence, as recorded in the arXiv licence metadata for this exact version and shown on the abstract page https://arxiv.org/abs/2608.10253v1. A full rights notice is delivered at licenses/arxiv-2608.10253-CC-BY-4.0.txt. Characters: every retained excerpt is pure ASCII, so the delivered engine, XeLaTeX with its default Latin Modern text font, reports no missing character.
\begin{lemma} \label{lemma:circles_canonical_embedding}
    Fix numbers $0 < a \le b < \infty$. Let $r : \mathbb{R}^2 \to (a, b)$ be a differentiable function for which $r_x \le 0$ and $r_y \ge 0$ and with bounded gradient. If $\alpha$ and $\beta$ are real numbers for which $\frac{a}{20} < \alpha - \beta < \frac{a}{10}$, then the associated canonical embedding $\mathsf{H}_{(\alpha, \beta)} = (\varphi_{\alpha}, \varphi_{\beta})$ is bilipschitz on the vertical strip with $\alpha < x < \alpha + \frac a {10}.$ The bilipschitz constant for $\mathsf{H}_{(\alpha, \beta)}$ depends on $a, b$, and $r$ but not on the selection of $\alpha$ and $\beta$.
\end{lemma}

\begin{proof}
    We begin with a calculation. Since $r$ is differentiable on its domain, the gradient of $\varphi_{\alpha}$ can be calculated as
    $$\nabla \varphi_{\alpha}(x, y) = \begin{bmatrix} \frac{r(x, y) r_x(x, y) + (\alpha - x)}{\sqrt{r^2(x, y) - (\alpha - x)^2}} \\ 1 + \frac{r(x, y) r_y(x, y)}{\sqrt{r^2(x, y) - (\alpha - x)^2}}\end{bmatrix} =: \begin{bmatrix} f_{\alpha}(x, y) \\ 1 + g_{\alpha}(x, y)\end{bmatrix}.$$
    Fix $\alpha < \beta$ and form the associated canonical embedding
    $$\mathsf{H}_{(\alpha, \beta)}(z) = (\varphi_{\alpha}(z), \varphi_{\beta}(z)).$$
    Using the gradient calculation, we find that
    $$\det D\mathsf{H}_{(\alpha, \beta)}(z) = \underbrace{f_{\alpha}(z) - f_{\beta}(z)}_{(I)} + \underbrace{\big(f_{\alpha}(z) g_{\beta}(z) - f_{\beta}(z) g_{\alpha}(z)\big)}_{(II)}.$$
    Our goal is to show that this is bounded above and below on the appropriate domain. Suppressing the arguments of the function $r$ for ease of notation, we can rewrite term (II) as 
    \begin{align*}
        (II) &= \frac{\big(r \cdot r_x + (\alpha - x)\big)\big(r\cdot r_y\big) - \big(r \cdot r_x + (\beta - x)\big)\big(r\cdot r_y\big)}{\sqrt{r^2 - (\alpha - x)^2}\sqrt{r^2 - (\beta - x)^2}} \\
        &= \frac{r \cdot r_y \cdot \big((\alpha - x) - (\beta - x)\big)}{\sqrt{r^2 - (\alpha - x)^2}\sqrt{r^2 - (\beta - x)^2}} \\
        &= \frac{r \cdot r_y \cdot (\alpha - \beta)}{\sqrt{r^2 - (\alpha - x)^2}\sqrt{r^2 - (\beta - x)^2}}.
    \end{align*}
    Likewise, term (I) can be rewritten as
    \begin{align*}
        (I) &= r \cdot r_x \left[\frac{1}{\sqrt{r^2 - (x - \alpha)^2}} - \frac{1}{\sqrt{r^2 - (x - \beta)^2}}\right] + \left[\frac{x - \beta}{\sqrt{r^2 - (x - \beta)^2}} - \frac{x - \alpha}{\sqrt{r^2 - (x - \alpha)^2}}\right].
    \end{align*}
    Combining (I) and (II) and refactoring, we have
    \begin{align*}
        \det D\mathsf{H}_{(\alpha, \beta)}(x, y) &= \left[\frac{(x - \beta) / r}{\sqrt{1 - ((x - \beta)/r)^2}} - \frac{(x - \alpha) / r}{\sqrt{1 - ((x - \alpha)/r)^2}} \right] \\
        &\quad\quad + r_x \left[\frac{1}{\sqrt{1 - ((x - \alpha) / r)^2}} - \frac{1}{\sqrt{1 - ((x - \beta)/r)^2}}\right] \\
        &\quad\quad + \frac{r_y}{r} \left[\frac{\alpha - \beta}{\sqrt{1 - ((x - \alpha)/r)^2} \sqrt{1 - ((x - \beta)/r)^2}}\right] \\
        &= (\text{A}) + (\text B) + (\text C).
    \end{align*}

For term (A), note that the function $\psi(t) = \frac{t}{\sqrt{1 - t^2}}$ is strictly increasing, and in fact $\psi'(t) \ge 1$. As a consequence, we have that
$$(\text{A}) \ge \frac{x - \beta}{r} - \frac{x - \alpha}{r} = \frac{\alpha - \beta}{r} \ge \frac{\alpha - \beta}{\|r\|_{\infty}} \ge \frac{a/20}{b} >  0.$$
Term (A) is also bounded above, since both $(x - \beta)/r$ and $(x - \alpha)/r$ lie in the interval $(0, 1/2)$. Note that these bounds depend on $r$ (and hence on $a$ and $b$), but not on $\alpha$ or $\beta$. 

Term (B) can be handled in a similar manner. Note that $\eta(t) = \frac{1}{\sqrt{1 - t^2}}$ is a strictly increasing function, so the bracketed term is negative (but as in Term (A), it is bounded in absolute value). Because $r_x \le 0$, it follows that $(\text B) \ge 0$ but is bounded above.

Term (C) can also be easily bounded: because $(x - \beta)/r$ and $(x - \alpha)/r$ both lie in $(0, 1/2)$, the bracketed portion is between $\alpha - \beta$ and $\frac 4 3 (\alpha - \beta)$. In any case, since $r_y/r$ is nonnegative and bounded above, the overall term is nonnegative and bounded. 

Combining the three upper and lower bounds, we conclude that the Jacobian determinant is bounded away from $0$ and $\infty$ with constants independent of $\alpha$ and $\beta$, as claimed. 
\end{proof}

\end{document}
